\begin{frame}{El gasto tributario llegó a 8,9\% del PIB en el año gravable 2024}
\centering
\small
\begin{tabular}{lrrrr}
\toprule
\textbf{Impuesto} & \textbf{2021} & \textbf{2022} & \textbf{2023} & \textbf{2024} \\
\midrule
IVA                        & 5,6 & 5,3 & 5,6 & 5,8 \\
Renta de personas naturales & 1,5 & 1,6 & 1,6 & 1,7 \\
Renta de personas jurídicas & 1,6 & 1,9 & 1,3 & 1,5 \\
\midrule
\textbf{Total}             & \textbf{8,7} & \textbf{8,8} & \textbf{8,5} & \textbf{8,9} \\
\bottomrule
\end{tabular}
\vspace{0.25cm}
\scriptsize El MFMP 2026 reporta esta estimación hasta el año gravable 2024; no publica todavía un valor para 2025.
\note{[Sources] MFMP 2026, apéndice, tabla A.2.1.}
\end{frame}

\begin{frame}{El gasto tributario aumentó 0,4 puntos del PIB entre 2023 y 2024}
\centering
\small
\begin{tabular}{lrrr}
\toprule
\textbf{Impuesto} & \textbf{2023} & \textbf{2024} & \textbf{Cambio} \\
\midrule
IVA                        & 5,6 & 5,8 & 0,2 \\
Renta de personas naturales & 1,6 & 1,7 & 0,1 \\
Renta de personas jurídicas & 1,3 & 1,5 & 0,2 \\
\midrule
\textbf{Total}             & \textbf{8,5} & \textbf{8,9} & \textbf{0,4} \\
\bottomrule
\end{tabular}
\vspace{0.2cm}
\scriptsize Las diferencias pueden no sumar por redondeo.
\note{[Sources] MFMP 2026, apéndice, tabla A.2.1.}
\end{frame}

\begin{frame}{Los bienes y servicios excluidos explican la mayor parte del gasto tributario del IVA}
\centering
\small
\begin{tabular}{lr}
\toprule
\textbf{Categoría, año gravable 2024} & \textbf{\% del PIB} \\
\midrule
Bienes y servicios excluidos          &  4,6 \\
Bienes y servicios exentos            &  1,2 \\
Tarifa reducida del 5\%               &  0,4 \\
No deducibilidad del IVA en bienes de inversión & -0,4 \\
\midrule
\textbf{Total}                        & \textbf{5,8} \\
\bottomrule
\end{tabular}
\note{[Sources] MFMP 2026, apéndice, tabla A.2.1.}
\end{frame}